\selectlanguage{ngerman}%
\frenchspacing%
\begin{ledgroupsized}[r]{120mm}%
\footnotesize%
\pstart%
\noindent\textbf{Überlieferung:}%
\pend%
\end{ledgroupsized}%
\begin{ledgroupsized}[r]{114mm}%
\footnotesize%
\pstart%
\parindent -6mm%
\makebox[6mm][l]{\textit{L}}%
Aufzeichnung:
LH XXXV 14,\;2 Bl.\;45.
Ein Blatt 4\textsuperscript{o},
Ränder zumeist beschnitten;
römisches\indexO{Rom} Papier (ohne Wasserzeichen).
Zwei ganzseitig beschriebene Seiten. %
\pend%
\end{ledgroupsized}%
\vspace{5mm}%
\begin{ledgroup}%
\footnotesize%
\pstart%
\noindent%
\textbf{Datierungsgründe:}
Die Datierung der Aufzeichnung N.\;\ref{58240} beruht auf Eigenschaften des Textträgers. %
Siehe die ausführliche Begründung auf S.\;\pageref{58240_dat}.%
\pend%
\end{ledgroup}%
\selectlanguage{latin}%
\frenchspacing%
\vspace{8mm}%
\normalsize%
\pstart%
\noindent%
\bracks{%
\href{\Scan{LH_35_14_02_045r}}%
{45\,r\textsuperscript{o}}%
}
Videndum jam,
an non sint lineae%
\indexS{centrum gravitatis linearum}
quarum centrum relativum%
\indexS{centrum gravitatis relativum}
\edtext{%
\bracks{vel}%
}{%
\lemma{vel}%
\Bfootnote{%
\textit{erg. Hrsg.}%
}}
absolutum%
\indexS{centrum gravitatis absolutum}
inveniri possit,
et quae proinde aptiores sint resolutioni figurae,
quam rectae.
Videamus nempe de circulo ejusque arcu.%
\indexS{circulus}%
\indexS{arcus circuli}
Sit ergo centrum arcus%
\indexS{centrum arcus circuli}
\edtext{%
circuli \textit{A},
centrum gravium \textit{E},%
\indexS{centrum gravium}%
\indexS{grave}%
}{%
\lemma{circuli \textit{A},}%
\Bfootnote{%
\textit{(1)}~punctum %
\textit{(2)}~centrum gravium \textit{E},%
~\textit{L}%
}}
ex quo quaeratur arcus dati momentum.%
\indexS{momentum arcus circuli}
Sit $ AB = x $
et $ AL = a $
et $ BL = \sqrt{aa-xx} $
et $ EA =
\edtext{%
[b] $%
}{%
\lemma{\textit{a}}%
\Bfootnote{%
\textit{L\;ändert Hrsg.}%
}}
et $ EB = b-x $,
fiet $ LE = \sqrt{bb-2bx\;\ovalbox{+xx}+aa\;\ovalbox{-xx}}$.
Sed haec ducenda in elementum arcus%
\indexS{elementum arcus circuli}
\edtext{%
circuli,
et fiet
$\displaystyle a \!\!\int\!\! \overline{dx\sqrt{\overline{bb+aa-2bx}:\overline{aa-xx}}}$.%
\rule[-5.6mm]{0pt}{12.4mm}%
}{%
\lemma{circuli,}%
\Bfootnote{%
\bars{quod est \textit{streicht Hrsg.}} %
\textit{(1)}~\bars{\textit{a} \textit{erg.}} $d \overline{x}:\sqrt{aa-xx}$ %
\textit{(2)}~et fiet $\displaystyle a \!\int\!\! \overline{dx \sqrt{\overline{bb+aa-2bx}:\overline{aa-xx}}}$.%
~\textit{L}%
}}
Hujus figurae datur area,%
\indexS{figura}%
\indexS{area figurae}
\edtext{%
opinor.
Saltem enim in eo casu dabitur,
quo arcus est integra circumferentia,%
\indexS{arcus circuli}%
\indexS{circumferentia integra}
ubi coincidere debere videtur absolutum centrum grav.%
\indexS{centrum gravitatis absolutum}
cum relativo.%
\indexS{centrum gravitatis relativum}%
}{%
\lemma{opinor.}%
\Bfootnote{%
\textit{(1)}~Sit enim %
\textit{(a)}~$\sqrt{b}$ %
\textit{(b)}~$x=ge+h$, fiet $\displaystyle \!\int\!\! \overline{de \sqrt{\overline{bb+aa-2bh-2be} : \overline{aa-hh-2hge-ggee}}}$. Fingatur $\sqrt{\vphantom{l}m+nx-ggee}$ ejus differentia est \textit{n} %
\textit{(2)}~Datur enim moment %
\textit{(3)}~Saltem enim %
\bracks{...} casu dabitur, %
\textit{(a)}~quo recta \textit{AE} bisecat arcum, tunc enim idem opinor %
\textit{(b)}~quo arcus %
\bracks{...} cum relativo.%
~\textit{L}%
}}
\pend%
\begin{edfigure}[2.25em][0.00em]
\includegraphics[scale=1]{%
gesamttex/edit_VIII-4/images/LH_35_14_02_045_d1_045r.pdf}
\edcaption[0.00em]{%
\bracks{\textit{Fig.\:1}}%
}
\end{edfigure}
\newpage%
\pstart%
Est et specialis difficultas%
\indexS{difficultas specialis}
in curvis,%
\indexS{curva}
quod etsi detur momentum%
\indexS{momentum curvae}
ex puncto \textit{E},
tamen non ideo datur centrum gravitatis,%
\indexS{centrum gravitatis curvae}
neque enim id semper cadit in ipsam lineam.
Ideo non sufficit distantiam ejus ab \textit{E} dari
nisi tunc,
cum
\edtext{%
recta
ut \textit{EB} lineam propositam%
}{%
\lemma{recta}%
\Bfootnote{%
\bars{ut \textit{erg.}} \textit{EB} %
\textit{(1)}~arcum %
\textit{(2)}~lineam propositam%
~\textit{L}%
}}
in duas partes congruas dividit.%
\indexS{linea divisa in duas partes congruas}
Poterimus tamen
ni fallor
invenire tali artificio,%
\indexS{artificium}
ut
\edtext{%
data linea
cujus centrum relativum quaeritur,%
\indexS{centrum gravitatis linearum}%
\indexS{centrum gravitatis relativum}
datoque puncto
quod est centrum gravium,%
\indexS{centrum gravium}%
\indexS{grave}
investigemus lineam
qua
\bracks{linea data} %
accedit ad centrum gravium%
\indexS{centrum gravium}%
\indexS{grave}%
\bracks{;}%
}{%
\lemma{data}%
\Bfootnote{%
\textit{(1)}~linea inquiramus in qua %
\textit{(2)}~linea cujus %
\textit{(a)}~momentum %
\textit{(b)}~centrum relativum quaeritur %
\textit{(aa)}~dataque %
\textit{(bb)}~datoque puncto %
\bracks{...} lineam qua %
\textit{(aaa)}~grave %
\textit{(bbb)}~\bars{linea data \textit{erg. Hrsg.}} accedit ad %
\textit{(aaaa)}~centrum %
\textit{(bbbb)}~centrum gravium%
~\textit{L}%
}}
seu dato ejus uno situ,%
\indexS{situs (corporis) gravis}
quaeramus situm proximum,
quo omnium %
maxime in dato tempore accedat,
et quaeramus
quodnam tunc punctum
(sive in linea,
sive
\edtext{%
extra figuram%
}{%
\lemma{extra}%
\Bfootnote{%
\textit{(1)}~figura
\textit{(2)}~figuram%
~\textit{L}%
}}
in superficie%
\indexS{superficies mota}
quam cum linea motam supponimus,
ita tamen
ut totius lineae pondus consideretur)%
\indexS{pondus lineae}
recta accesserit ad centrum gravium%
\indexS{centrum gravium}%
\indexS{grave}%
\bracks{:}
id punctum erit centrum gravitatis relativum.%
\indexS{centrum gravitatis relativum}
Hoc subtilissimum,
sed opinor satis
\edlabel{LH_35_14_02_045_045r_sxn-1}%
difficile.%
\edtext{}{%
\xxref{LH_35_14_02_045_045r_sxn-1}{LH_35_14_02_045_045r_sxn-2}%
\lemma{difficile}%
\Bfootnote{%
\textit{(1)}~; est et alia via brevior pro arcu circuli %
\textit{(2)}~. \bracks{/} Si recta rigida%
~\textit{L}%
}}%
\pend%
\begin{edfigure}[1.0em][1.0em]
\includegraphics[scale=1]{%
gesamttex/edit_VIII-4/images/LH_35_14_02_045_d2_045r.pdf}
\edcaption{\lbrack\textit{Fig.\:2}\rbrack}
\end{edfigure}
\pstart%
Si recta rigida%
\edlabel{LH_35_14_02_045_045r_sxn-2}%
\indexS{recta rigida}
ducatur ex centro gravium%
\indexS{centrum gravium}%
\indexS{grave}
ad centrum gravitatis,
impediet
\edtext{%
appropinquationem gravis,%
\indexS{appropinquatio gravis}%
}{%
\lemma{appropinquationem}%
\Bfootnote{%
\bars{puncti \textit{gestr.}} gravis,%
~\textit{L}%
}}
et hoc est quasi aequilibrium in linea recta.%
\indexS{aequilibrium in linea recta}
An vero hinc sequitur
nec ipsum corpus circa hoc punctum converti?
Ita videtur;
neque enim motus erit,%
\indexS{motus gravis}
sine aliqua accessione.%
\indexS{accessio gravis}
Quod si duae rectae rigidae%
\indexS{recta rigida}
et diversa corpora pugnent%
\indexS{corpora pugnantia}
ut \textit{A} et \textit{B} simul
\edtext{%
contra \textit{C},%
}{%
\lemma{contra}%
\Bfootnote{%
\bars{contra \textit{streicht Hrsg.}} \textit{C},%
~\textit{L}%
}}
et \textit{A} et \textit{C} contra \textit{B},
et \textit{B} et \textit{C} contra \textit{A}.
Quaerendum ergo centrum commune omnium%
\indexS{centrum gravitatis commune}
quod feretur in \textit{E}.
\bracks{%
\href{\Scan{LH_35_14_02_045v}}%
{45\,v\textsuperscript{o}}%
}
\pend%
\vspace{1.0em}%
\pstart%
\noindent
\bracks{\textit{Nachfolgend kleingedruckter Text in L gestrichen:}}
\pend%
\vspace{0.5em}%
\footnotesize%
\pstart%
Tractatus aliquis de hac materia%
\indexS{materia tractatus}
posset appellari Geocentricus%
\indexS{tractatus geocentricus}
\edtext{%
cum quaeritur%
}{%
\lemma{cum}%
\Bfootnote{%
\textit{(1)}~supponitur %
\textit{(2)}~quaeritur%
~\textit{L}%
}}
quid gravia factura sint
in locis centro gravium vicinis.%
\indexS{grave}%
\indexS{centrum gravium}%
\indexS{locus vicinus centro gravium}
\pend%
\pstart%
Ex punctis quotcunque quadrato ordine%
\indexS{ordo quadratus}
\normalsize{\bracks{\textit{Text bricht ab.}}}
\pend%
\normalsize%
\begin{edfigure}[0.00em][2.50em]%
\hspace{-85mm}%
\includegraphics[scale=1]{%
gesamttex/edit_VIII-4/images/LH_35_14_02_045_d3_045v.pdf}
\edcaption{%
\hspace{-85mm}%
\bracks{\textit{Fig.\:3, gestr.}}%
}%
\vspace{-6.70em}%
\hspace{1mm}%
\includegraphics[scale=1]{%
gesamttex/edit_VIII-4/images/LH_35_14_02_045_d4_045v.pdf}
\edcaption{%
\hspace{1mm}%
\bracks{\textit{Fig.\:4, gestr.}}%
}%
\vspace{-6.70em}%
\hspace{85mm}%
\includegraphics[scale=1]{%
gesamttex/edit_VIII-4/images/LH_35_14_02_045_d5_045v.pdf}
\edcaption{%
\hspace{85mm}%
\bracks{\textit{Fig.\:5, gestr.}}%
}%
\vspace{2.0em}%
\hspace{-45mm}%
\includegraphics[scale=1]{%
gesamttex/edit_VIII-4/images/LH_35_14_02_045_d6_045v.pdf}
\edcaption{%
\hspace{-45mm}%
\bracks{\textit{Fig.\:6, gestr.}}%
}
\vspace{-7.5em}%
\hspace{55mm}%
\includegraphics[scale=1]{%
gesamttex/edit_VIII-4/images/LH_35_14_02_045_d7_045v.pdf}
\edcaption{%
\hspace{50mm}%
\bracks{\textit{Fig.\:7, gestr.}}%
}%
\end{edfigure}
\pstart%
Sit
\edtext{%
recta quaecunque \textit{BC};%
}{%
\lemma{recta quaecunque \textit{BC}}%
\Cfootnote{%
Vgl. das Diagramm \bracks{\textit{Fig.\:8}}.%
}}
ducatur in distantias a
\edtext{%
puncto dato
\bracks{\textit{E}}.%
}{%
\lemma{puncto}%
\Bfootnote{%
\textit{(1)}~suo \textit{C} %
\textit{(2)}~dato %
\bars{\textit{E} \textit{erg. Hrsg. nach \textup{[}Fig.\:8\textup{]}}}.%
~\textit{L}%
}}
Sit $ BV = v $,
quaeritur \textit{VE}.
Sit $ BF = b $,
et $ EF = c $,
fiet
$ \sqrt{\vphantom{l^i}cc+bb \MP 2bv+vv} = EV $.%
\rule[-2.0mm]{0pt}{7.5mm}
Ergo summa omnium \textit{EV}
insistentium in \textit{V},
erit
\edtext{%
$ \displaystyle\!\!\int\!\! \overline{dv\sqrt{cc+bb \MP 2bv+vv}} $.%
\rule[-2.5mm]{0pt}{8.0mm}
Haec summa pendet ex quadratura Hyperbolae.%
\indexS{quadratura hyperbolae}%
\indexS{hyperbola}%
}{%
\lemma{$ \displaystyle\!\int\!\! \overline{dv\sqrt{cc+bb \MP 2bv+vv}} $}%
\Bfootnote{%
\bars{$=$ . Quae pendet ex quadratura Hyperbolae \textit{streicht Hrsg.}}.
\textit{(1)}~Si \textit{FV} vocetur \textit{e}, fiet $ \displaystyle\!\!\int\!\! \overline{\sqrt{cc+ee}\;d\overline{e}}$ %
\textit{(2)}~Haec summa %
\bracks{...} quadratura Hyperbolae.%
~\textit{L}%
}}
Porro maxima \textit{v} est \textit{y} seu \textit{BC}.%
\rule[-0.0mm]{0pt}{5.0mm}
Inventa igitur summa dicta,
pro \textit{v} substituatur \textit{y} seu \textit{BC},
et pro \textit{y} ejus valor%
\indexS{valor}
\edtext{%
ex \textit{AB}%
}{%
\lemma{ex}%
\Bfootnote{%
\textit{(1)}~\textit{AE} %
\textit{(2)}~\textit{AB}%
~\textit{L}%
}}
seu \textit{x},
et fiet:
$ \displaystyle\!\!\int\!\! \overline{\displaystyle\!\!\int\!\! \overline{dv\sqrt{cc+bb \MP 2bv+vv}}\;dx}$.%
\rule[-4.0mm]{0pt}{10.5mm}%
\pend%
\pstart%
Caeterum ex his patet
regulam centri gravitatis%
\indexS{regula centri gravitatis}
hic non succedere,
qua
\edtext{%
alioqui%
}{%
\lemma{alioqui}%
\Cfootnote{%
Nicht ermittelt.%
}}
utimur,
ut scilicet habeatur
exempli causa
momentum rectae \textit{BC}%
\indexS{momentum rectae}
ex puncto \textit{E},
\edtext{%
ducendo rectam \textit{BC}%
}{%
\lemma{ducendo}%
\Bfootnote{%
\textit{(1)}~punctum %
\textit{(2)}~rectam \textit{BC}%
~\textit{L}%
}}
in \textit{EV} distantiam centri gravitatis%
\indexS{distantia centri gravitatis}
seu medii ejus,
quod est centrum gravitatis absolutum.%
\indexS{centrum gravitatis absolutum}
Et proinde sequi dari quoddam centrum
\edtext{%
gravitatis \textit{G}
rectae ut \textit{BC},
relativum%
\indexS{centrum gravitatis relativum}%
}{%
\lemma{gravitatis}%
\Bfootnote{%
\bars{\textit{G} \textit{erg.}} %
\textit{(1)}~relativum %
\textit{(2)}~rectae ut \textit{BC}, relativum%
~\textit{L}%
}}
ad centrum gravium \textit{E}.%
\indexS{centrum gravium}%
\indexS{grave}
Nempe momentum rectae%
\indexS{momentum rectae}
\edtext{%
\textit{BC}%
}{%
\lemma{\textit{BC}}%
\Bfootnote{\textit{erg.\;L}%
}}
est
$ \!\displaystyle\!\int\!\! \overline{dv\sqrt{cc+bb \MP 2bv+vv}} $,%
\rule[-3.0mm]{0pt}{8.5mm}
posito maximum\;\textit{v} esse ipsam rectam \textit{BC};
\ledpb%
jam dividatur hoc momentum%
\indexS{momentum rectae}
per \textit{BC} seu \textit{y}
(seu applicetur ipsi maxima \textit{v})
fiet
$ \!\displaystyle\!\int\!\! dv\;\overline{\sqrt{cc+bb \MP 2bv+vv}} : y = EG $.
Seu centro \textit{E},
radio%
\indexS{radius circuli}
\edtext{%
\bracks{seu}%
}{%
\lemma{seu}%
\Bfootnote{%
\textit{erg. Hrsg.}%
}}
valore isto ipsius \textit{EG}%
\indexS{valor}
describatur circulus,%
\indexS{circulus}%
\rule[-3.0mm]{0pt}{6.6mm}
secans rectam \textit{BC} %
in puncto \textit{G} quaesito,
quod erit ejus centrum gravitatis relativum.%
\indexS{centrum gravitatis relativum}
\edtext{%
Etenim hoc supponit quadraturam Hyperbolae,%
\indexS{quadratura hyperbolae}%
\indexS{hyperbola}
non potest\;\textit{G} cadere regulariter in \angles{ra}dium.%
\indexS{radius circuli}%
}{%
\lemma{Etenim hoc}%
\Bfootnote{%
\bracks{...} in \angles{ra}dium. %
\textit{erg.\;L}%
}}
Hinc si \textit{G} cadat in \textit{F}
seu rectam \textit{BC} continuatam,
tunc centrum gravitatis absolutum%
\indexS{centrum gravitatis absolutum}
coincidet cum relativo,%
\indexS{centrum gravitatis relativum}
nam
\edtext{%
posita \textit{FE} seu $c=0$,
fiet%
}{%
\lemma{posita}%
\Bfootnote{%
\textit{(1)}~$G=0$ fi %
\textit{(2)}~$c=0$ se %
\textit{(3)}~\textit{FE} seu $c=0$, fiet%
~\textit{L}%
}}
\textit{EG} seu
$ FG = \displaystyle\frac{1}{y} \!\int\!\! \overline{dv\;\overline{b-v}} = \overline{bv - \displaystyle\frac{1}{2}vv} : y$. %
Seu quia maxima $v=y$,
fiet $FG=b-\displaystyle\frac{1}{2}v$.
Ergo $BG=\displaystyle\frac{1}{2}v$, %
\edtext{%
quando \textit{E} coincidit
\bracks{cum}
\textit{F}.%
}{%
\lemma{quando}%
\Bfootnote{%
\textit{(1)}~\textit{C} %
\textit{(2)}~\textit{E} coincidit %
\bars{cum \textit{erg. Hrsg.}}%
~\textit{L}%
}}
\pend%
\begin{edfigure}[1.25em][2.0em]
\includegraphics[scale=1]{%
gesamttex/edit_VIII-4/images/LH_35_14_02_045_d8_045v.pdf}
\edcaption[0.50em]{%
\bracks{\textit{Fig.\:8}}%
}%
\end{edfigure}%
\pstart%
\noindent%
\bracks{\textit{Auf Bl.\;\href{\Scan{LH_35_14_02_045r}}{45\,r\textsuperscript{o}} schräg zur Schreibrichtung:}}
\pend%
\vspace{0.5em}%
\pstart%
\noindent%
Per centr.\;gr. reducuntur corpora%
\indexS{corpus reductum ad centrum gravitatis}
\edtext{%
ad puncta%
\bracks{,}
puncta in aequilibrio sunt%
\indexS{aequilibrium punctorum}
si sint reciproce%
}{%
\lemma{ad puncta}%
\Bfootnote{%
\textit{(1)}~corp %
\textit{(2)}~puncta %
\textit{(a)}~ponderant recipr %
\textit{(b)}~in aequilibrio %
\bracks{...} sint reciproce%
~\textit{L}%
}}
ut distantiae circulares%
\indexS{distantia circularis}
seu in sphaera%
\indexS{sphaera}
circa centrum descripta,
minima per puncta duo et punctum suspensionis%
\indexS{punctum suspensionis librae}%
\indexS{suspensio librae}
\edtext{%
libr\bracks{a}e%
\indexS{libra}%
}{%
\lemma{libre}%
\Bfootnote{%
\textit{L\;ändert Hrsg.}%
}}
ducta.
Sed quid
si haec tria puncta una sint
in eodem centro maximo?%
\indexS{centrum maximum}
\pend%
\pstart%
Videndum etiam
an non centrum gravitatis%
\indexS{centrum gravitatis relativum}
\edtext{%
relativum%
}{%
\lemma{relativum}%
\Bfootnote{%
\textit{erg.\;L}%
}}
durante motu
mutetur in corpore.
\pend%
\pstart%
Videtur corpus
ex centro gravitatis suspensum%
\indexS{corpus suspensum e centro gravitatis}
eodem modo se habere,
ac si ejus vis in centrum illud reducta esset%
\indexS{vis corporis reducta in centrum gravitatis}
cum durante motu mutetur centrum grav.%
\indexS{centrum gravitatis mutans durante motu}
\pend%
\pstart%
Videndum an nihilominus verum sit
viam centri gravitatis,%
\indexS{via centri gravitatis}
seu potius diversorum compositam
in unum ductam%
\bracks{,}
in corpus aequari areae.
\pend%
\newpage%
\pstart%
\noindent%
\centering%
*     *     *
\pend%
\vspace{4.0em}%
\selectlanguage{german}%
\frenchspacing%
\footnotesize%
\pstart%
\noindent%
\textls{Ausführliche Datierungsgründe:}%
\pagelabel{58240_dat}
Die titellose und undatierte Aufzeichnung N.\;\ref{58240} ist auf einem Träger überliefert, der zwar keine Spuren eines Wasserzeichens aufweist, anhand seiner äußeren Merkmale aber als römisches\indexO{Rom} Papier identifiziert werden kann.
In der Farbe und der Anordnung der Stege ähnelt er nämlich den Trägern weiterer im vorliegenden Band edierter Texte (etwa N.\;\ref{58656-A}, N.\;\ref{58258}, N.\;\ref{58218}, N.\;\ref{60197}, N.\;\ref{58220} und N.\;\ref{58235}), deren römische\indexO{Rom} Herkunft feststellbar ist.
Auf gleichem Papier sind ebenfalls Leibnizens Auszüge mit Bemerkungen aus \textsc{P.\;de Fermat}, \textit{Varia opera mathematica} (Toulouse 1679)\cite{00533} überliefert, welche zum Teil unter N.\;\ref{58229} im vorliegenden Band ediert sind.
Auch die Aufzeichnung N.\;\ref{58240} ist demgemäß -- wie die Auszüge mit Bemerkungen N.\;\ref{58229} -- auf Leibnizens Aufenthalt in Rom\indexO{Rom} zu datieren, d.\,h. auf die Monate von Mitte April bis November 1689.
Eine frühere Entstehung ist ausgeschlossen, eine spätere unwahrscheinlich.
\pend%
\selectlanguage{latin}%
\frenchspacing%
\normalsize%
